% ===========================================================================
\section{Extensiones por logro}
\label{sec:logros}
% ===========================================================================

Se persiguen \textbf{los nueve logros de la Etapa 2 que afectan al analizador}, más el
\emph{carry-over} de los floats de la Etapa 1. Cada uno se transforma con el mismo procedimiento
de cuatro pasos; la diferencia con el núcleo es que acá los conflictos \textbf{no vienen de la
gramática de la cátedra: los crea el logro}.

\subsection{La condición de atomicidad}
\label{sec:atomicidad}

El compilador debe poder ejecutarse y probarse con cualquier subconjunto de logros, incluido el
vacío. A nivel de gramática eso se traduce en una condición dura:

\begin{quote}
\textbf{Para todo subconjunto de logros, la gramática resultante debe ser LL(1).}
\end{quote}

No alcanza con que lo sea con todos prendidos. La condición tiene una consecuencia de diseño
concreta: \textbf{ningún logro puede apoyarse en un no terminal que introdujo otro}. Cuando dos
logros necesitan una construcción parecida, cada uno lleva la suya --- ver
sección~\ref{sec:logro-inicializaciones}.

Como técnica, conviene que \textbf{desactivar un logro sea quitar alternativas y no cambiar de
estructura}. Los logros de ternario y de operadores posfijos son el ejemplo limpio: apagados,
sus no terminales quedan en sólo \eps{} y la capa desaparece sin tocar ninguna otra regla.

\subsection{Visibilidad Mejorada E2}
\label{sec:logro-visibilidad}

\paso{Defecto que crea} \tl{public} y \tl{private} pasan a poder encabezar cualquier miembro. El
\tl{public} del \nt{Constructor} original queda absorbido, el constructor se vuelve \tl{idClase}
\nt{CuerpoDeMetodo} y \textbf{empieza a competir} con el atributo o método de tipo referencia por
el token \tl{idClase}: \texttt{C()\{\}} contra \texttt{C x;}.

\paso{Transformación} Factorizar la visibilidad al frente y, adentro, factorizar el \tl{idClase}.
El desempate es el token que sigue: \tl{(} es constructor, cualquier otro abre un tipo.
\begin{producciones}
\regla{Miembro}{\nt{VisibilidadOpcional} \nt{MiembroSinVisibilidad}}
\regla{VisibilidadOpcional}{\tl{public} \alt{} \tl{private} \alt{} \eps}
\regla{MiembroSinVisibilidad}{\tl{static} \nt{TipoMetodo} \tl{idMetVar} \nt{CuerpoDeMetodo}}
\regla{MiembroSinVisibilidad}{\tl{void} \tl{idMetVar} \nt{CuerpoDeMetodo}}
\regla{MiembroSinVisibilidad}{\nt{TipoPrimitivo} \nt{DimensionesOpcionales} \tl{idMetVar} \nt{AtributoOMetodo}}
\regla{MiembroSinVisibilidad}{\tl{idGen} \nt{DimensionesOpcionales} \tl{idMetVar} \nt{AtributoOMetodo}}
\regla{MiembroSinVisibilidad}{\tl{idClase} \nt{MiembroDesdeIdClase}}
\regla{MiembroDesdeIdClase}{\nt{CuerpoDeMetodo}}
\regla{MiembroDesdeIdClase}{\nt{TipoGenericoOpcional} \nt{DimensionesOpcionales} \tl{idMetVar} \nt{AtributoOMetodo}}
\end{producciones}
$\{$\tl{static}$\}$, $\{$\tl{void}$\}$, $\{$\tl{boolean}, \tl{char}, \tl{int}, \tl{float}$\}$,
$\{$\tl{idGen}$\}$, $\{$\tl{idClase}$\}$: disjuntos. Y en \nt{MiembroDesdeIdClase}, $\{$\tl{(}$\}$
contra $\{$\tl{<}, \tl{[}, \tl{idMetVar}$\}$: disjuntos.

Con el logro apagado vuelve la forma del núcleo y \nt{MiembroDesdeIdClase} se queda con una sola
producción.

\begin{decision}{los métodos de interfaz no llevan visibilidad}
El enunciado habla de «Métodos, Atributos y Constructores», que son los \nt{Miembro} de una clase.
\nt{MetodoInterfaz} es un no terminal aparte y en Java los métodos de interfaz son implícitamente
públicos, así que se deja como está.
\end{decision}

\subsection{Atributos Inicializados y Variables locales Clásicas E2}
\label{sec:logro-inicializaciones}

Los dos agregan un \tl{=} \nt{ExpresionCompuesta} opcional, y la tentación es usar un único no
terminal compartido. \textbf{Está mal} por la condición de atomicidad: son logros distintos, y con
\emph{atrinit} apagado y \emph{varclasica} prendido las variables locales tienen que poder
inicializarse igual. Van dos no terminales dedicados, \nt{InicializacionDeAtributo} e
\nt{InicializacionDeDeclarador}, con el mismo cuerpo y distinta marca de logro.

\paso{Defecto que crea \emph{varclasica}} Una declaración local pasa a poder empezar con
\tl{idClase} (\texttt{C x;}), y una sentencia-expresión también (\texttt{C.m();}, por
\nt{LlamadaMetodoEstatico}). Chocan.

\paso{Transformación} Factorizar el \tl{idClase} en \nt{Sentencia}: \tl{.} significa expresión,
cualquier otro token significa declaración.
\begin{producciones}
\regla{Sentencia}{\nt{ExpresionSinIdClase} \tl{;}}
\regla{Sentencia}{\nt{TipoPrimitivo} \nt{DimensionesOpcionales} \nt{ListaDeclaradores} \tl{;}}
\regla{Sentencia}{\tl{idGen} \nt{DimensionesOpcionales} \nt{ListaDeclaradores} \tl{;}}
\regla{Sentencia}{\tl{idClase} \nt{SentenciaDesdeIdClase}}
\regla{SentenciaDesdeIdClase}{\nt{TipoGenericoOpcional} \nt{DimensionesOpcionales} \nt{ListaDeclaradores} \tl{;}}
\regla{SentenciaDesdeIdClase}{\tl{.} \tl{idMetVar} \nt{ArgsActuales} \nt{ColaDeExpresionDesdePrimario} \tl{;}}
\end{producciones}

La rama de expresión tiene que retomar una expresión \emph{ya empezada}. Escribirla a mano
duplicaría media jerarquía de expresiones; en su lugar se insertan tres no terminales que capturan
las colas compartidas:
\begin{producciones}
\regla{ExpresionSinIdClase}{\nt{OperadorUnario} \nt{Operando} \nt{ColaTrasOperando}}
\regla{ExpresionSinIdClase}{\nt{Primitivo} \nt{ColaTrasOperando}}
\regla{ExpresionSinIdClase}{\nt{PrimarioSinIdClase} \nt{ColaDeExpresionDesdePrimario}}
\regla{ColaTrasOperando}{\nt{RestoExpresionCompuesta} \nt{RestoCondicional} \nt{RestoExpresion}}
\regla{ColaDeExpresionDesdePrimario}{\nt{RestoReferencia} \nt{PosfijoOpcional} \nt{ColaTrasOperando}}
\regla{Primario}{\nt{LlamadaMetodoEstatico} \alt{} \nt{PrimarioSinIdClase}}
\end{producciones}
\nt{LlamadaMetodoEstatico} es la única forma de \nt{Primario} que empieza con \tl{idClase}, lo que
hace que la partición sea exacta. \textbf{Sólo se duplican las cabezas}: todas las colas se
comparten con la jerarquía normal, de modo que no hay lógica repetida y un cambio futuro en las
expresiones no hay que espejarlo.

\begin{decision}{alternativa descartada: dos tokens de anticipación}
El mismo conflicto se resuelve dándole al analizador \emph{dos} tokens de anticipación sólo en este
punto --- mirar si después del \tl{idClase} viene un \tl{.} --- con cero no terminales nuevos. Se
descartó porque despega la implementación del esquema de traducción visto en clase, que es de un
único token, y ese esquema es la garantía de correctitud que se está siguiendo. El costo de la
solución LL(1) pura son tres no terminales sin lógica duplicada, que es barato.
\end{decision}

\subsection{Fors! E2}

Un \nt{For} nuevo como alternativa de \nt{Sentencia}. El encabezado tiene el mismo problema de
\tl{idClase} que la sentencia y se resuelve igual, reutilizando \nt{ExpresionSinIdClase} y
\nt{ColaDeExpresionDesdePrimario}. El desempate entre el \tl{for} estándar y el de iteradores es de
un solo token y está justo después del declarador: \tl{:} es de iteradores, \tl{=} o \tl{;} es
estándar.
\begin{producciones}
\regla{For}{\tl{for} \tl{(} \nt{EncabezadoFor} \tl{)} \nt{Sentencia}}
\regla{EncabezadoFor}{\tl{;} \nt{ExpresionOpcional} \tl{;} \nt{ExpresionOpcional}}
\regla{EncabezadoFor}{\tl{var} \tl{idMetVar} \nt{EncabezadoForTrasVar}}
\regla{EncabezadoFor}{\nt{TipoPrimitivo} \nt{DimensionesOpcionales} \tl{idMetVar} \nt{EncabezadoForTrasDeclarador}}
\regla{EncabezadoFor}{\tl{idGen} \nt{DimensionesOpcionales} \tl{idMetVar} \nt{EncabezadoForTrasDeclarador}}
\regla{EncabezadoFor}{\tl{idClase} \nt{EncabezadoForDesdeIdClase}}
\regla{EncabezadoFor}{\nt{ExpresionSinIdClase} \nt{ClausulasFinalesFor}}
\regla{EncabezadoForTrasDeclarador}{\tl{:} \nt{Expresion}}
\regla{EncabezadoForTrasDeclarador}{\nt{InicializacionDeDeclaradorFor} \nt{ClausulasFinalesFor}}
\regla{ClausulasFinalesFor}{\tl{;} \nt{ExpresionOpcional} \tl{;} \nt{ExpresionOpcional}}
\end{producciones}
\nt{ClausulasFinalesFor} con una sola \nt{ExpresionOpcional} por separador es la lectura literal de
«se restringe a una sola sentencia o expresión por separador».

\subsection{Operador Ternario E2}

Una capa nueva entre la asignación y los operadores binarios, con la precedencia de Java.
\begin{producciones}
\regla{Expresion}{\nt{ExpresionCondicional} \nt{RestoExpresion}}
\regla{RestoExpresion}{\nt{OperadorAsignacion} \nt{ExpresionCondicional} \alt{} \eps}
\regla{ExpresionCondicional}{\nt{ExpresionCompuesta} \nt{RestoCondicional}}
\regla{RestoCondicional}{\tl{?} \nt{ExpresionCondicional} \tl{:} \nt{ExpresionCondicional} \alt{} \eps}
\end{producciones}
Asociativo a \textbf{derecha} --- \texttt{a ? b : c ? d : e} es \texttt{a ? b : (c ? d : e)} ---
porque la rama del \tl{:} vuelve a \nt{ExpresionCondicional}, que puede tener su propio
\nt{RestoCondicional}. Con el logro apagado, \nt{RestoCondicional} queda en sólo \eps{} y
\nt{ExpresionCondicional} se vuelve un alias de \nt{ExpresionCompuesta}: la capa desaparece sin
dejar rastro. Nótese que \nt{RestoExpresion} conserva \nt{ExpresionCondicional} y no
\nt{Expresion} a la derecha, preservando la no asociatividad de la asignación.

Este logro requiere el token \tl{?}, que \textbf{no existe en el léxico de la Etapa 1}.

\subsection{Operadores Posifijos E2}

\begin{producciones}
\regla{Operando}{\nt{Primitivo} \alt{} \nt{Referencia} \nt{PosfijoOpcional}}
\regla{PosfijoOpcional}{\tl{++} \alt{} \tl{-{}-} \alt{} \eps}
\end{producciones}
El posfijo va sobre \nt{Referencia} y \textbf{no} sobre \nt{Operando}: \texttt{5++} no tiene
sentido. Misma técnica de desactivación que el ternario. Es el único logro que \textbf{no}
requiere tocar el léxico: \tl{++} y \tl{-{}-} ya son tokens desde la Etapa 1 y no tenían dónde
usarse.

\subsection{MiniLambdas E2}

\paso{Defecto que crea} El \tl{(} abre tanto una expresión parentizada como la lista de parámetros
de una lambda, y \tl{idMetVar} abre tanto un acceso a variable como una lambda de un parámetro.
A primera vista \texttt{(a, b, c) -> x} contra \texttt{(a + b)} exige escanear hasta el \tl{)}
para decidir, o sea anticipación ilimitada.

\paso{Análisis} \textbf{No la exige.} Los casos se separan solos:

\begin{center}
\begin{tabular}{@{}lll@{}}
\toprule
\textbf{Entrada} & \textbf{Decisión} & \textbf{Tokens} \\
\midrule
\tl{(} \tl{)}            & sólo puede ser lambda de cero parámetros & 1 (el \tl{)}) \\
\tl{(} $e$ \tl{,} \dots  & sólo puede ser lambda: no hay operador coma & 1 (la \tl{,}) \\
\tl{(} $e$ \tl{)}        & lo decide el token \emph{posterior} al \tl{)} & 1 (el \tl{->}) \\
\tl{idMetVar} \tl{->}    & lo decide el token posterior al identificador & 1 (el \tl{->}) \\
\bottomrule
\end{tabular}
\end{center}

La clave es que el \tl{->} viene \emph{inmediatamente} después del \tl{)}, así que en el momento
de decidir sigue siendo un token de anticipación.

\paso{Resultado}
\begin{producciones}
\regla{ParentizadaOLambda}{\tl{(} \nt{ContenidoDeParentesis}}
\regla{ContenidoDeParentesis}{\tl{)} \tl{->} \nt{CuerpoLambda} \alt{} \nt{Expresion} \nt{RestoDeParentesis}}
\regla{RestoDeParentesis}{\tl{)} \nt{FinDeParentizadaOLambda}}
\regla{RestoDeParentesis}{\tl{,} \nt{ListaParametrosLambda} \tl{)} \tl{->} \nt{CuerpoLambda}}
\regla{FinDeParentizadaOLambda}{\tl{->} \nt{CuerpoLambda} \alt{} \eps}
\regla{AccesoVarLlamadaOLambda}{\tl{idMetVar} \nt{AccesoLlamadaOLambda}}
\regla{AccesoLlamadaOLambda}{\nt{ArgsActuales} \alt{} \tl{->} \nt{CuerpoLambda} \alt{} \eps}
\regla{CuerpoLambda}{\nt{Expresion}}
\end{producciones}
\nt{CuerpoLambda} \deriva{} \nt{Expresion} es la lectura literal de «el cuerpo debe estar formado
por una única expresión, no se permiten cuerpos delimitados por llaves»: como no hay alternativa
con \tl{\{}, la gramática \textbf{rechaza} \texttt{x -> \{ \dots \}} sin ningún chequeo adicional.

Requiere el token \tl{->}, que no existe en el léxico de la Etapa 1.

\begin{decision}{permisividad declarada de las lambdas}
En \texttt{( a , b ) -> x} el \texttt{a} se deriva como \nt{Expresion}, no como parámetro, de modo
que \texttt{( a + 1 , b ) -> x} \textbf{atraviesa el análisis sintáctico}. Rechazarlo exige saber
que esa expresión tiene que ser un identificador suelto, lo que es un chequeo semántico. Es la
misma clase de permisividad que la gramática original ya tiene con \texttt{5 = x;}. Lo mismo vale,
como se anticipó en la sección~\ref{sec:alcance-chequeos}, para la prohibición de capturar
variables.
\end{decision}

\subsection{Inicializacion de Arreglos E2}

Después del \tl{[} de una creación de arreglo, o viene una expresión de tamaño o viene el \tl{]}
seguido del inicializador. Un token de desempate.
\begin{producciones}
\regla{CreacionDeArreglo}{\tl{[} \nt{RestoCreacionDeArreglo}}
\regla{RestoCreacionDeArreglo}{\nt{Expresion} \tl{]} \nt{DimensionesConTamanioOpcional}}
\regla{RestoCreacionDeArreglo}{\tl{]} \tl{\{} \nt{ListaExpsOpcional} \tl{\}}}
\end{producciones}
Se reutiliza \nt{ListaExpsOpcional}, que ya existe para los argumentos actuales, porque es la misma
construcción. La restricción «sólo se permite como parte de una expresión de creación utilizando
\tl{new}» queda impuesta por la gramática: \nt{CreacionDeArreglo} sólo es alcanzable desde
\nt{Creacion}.

\subsection{Genericidad Avanzada E2}

Dos cosas distintas. El \textbf{anidamiento} sale de una sola recursión:
\begin{producciones}
\regla{ArgumentoDeTipo}{\tl{idGen} \alt{} \tl{idClase} \nt{TipoGenericoOpcional}}
\end{producciones}
Con el logro apagado la segunda alternativa es \tl{idClase} a secas. (\texttt{List<List<Persona>>}
cierra con dos \tl{>} separados: el léxico no define ningún token \texttt{>>}, así que la regla de
\emph{maximal munch} entrega dos \tl{>} y no hay nada que resolver del lado sintáctico.)

El \textbf{diamante} exige más cuidado, porque el enunciado dice que «sólo puede utilizarse al
instanciar una Clase Genérica». Se impone \textbf{por gramática}, separando el genérico en posición
de tipo del genérico bajo \tl{new}:
\begin{producciones}
\regla{TipoGenericoOpcional}{\tl{<} \nt{ArgumentoDeTipo} \tl{>} \alt{} \eps}
\regla{CreacionDesdeIdClase}{\tl{<} \nt{GenericoDeCreacion} \alt{} \nt{CreacionDeArreglo} \alt{} \nt{ArgsActuales}}
\regla{GenericoDeCreacion}{\tl{>} \nt{ArgsActuales}}
\regla{GenericoDeCreacion}{\nt{ArgumentoDeTipo} \tl{>} \nt{ArregloOInstanciaDeClase}}
\end{producciones}
La primera alternativa de \nt{GenericoDeCreacion} \emph{es} el diamante, y va sólo hacia
\nt{ArgsActuales}. Resultado:

\begin{center}
\begin{tabular}{@{}ll@{}}
\toprule
\textbf{Entrada} & \textbf{Veredicto} \\
\midrule
\texttt{List<Persona> x;}        & acepta \\
\texttt{List<List<Persona>> x;}  & acepta \\
\texttt{new List<>()}            & acepta \\
\texttt{new List<Persona>[3]}    & acepta \\
\texttt{List<> x;}               & \textbf{rechaza} (diamante en posición de tipo) \\
\texttt{new List<>[3]}           & \textbf{rechaza} (diamante con arreglo) \\
\bottomrule
\end{tabular}
\end{center}

\begin{decision}{compartir un no terminal puede volverse incorrecto más adelante}
En la gramática original, \nt{TipoGenericoOpcional} se usa desde \nt{TipoReferencia} y desde
\nt{LlamadaConstructor}, y compartirlo es correcto: son la misma construcción. \textbf{Deja de
serlo con este logro}, porque el diamante vale en una posición y no en la otra. Compartiéndolo
habría que aceptar \texttt{List<> x;} y rechazarlo después en el análisis semántico. Separarlos
cuesta un no terminal y mueve la restricción al lugar donde se puede verificar sin conocer el
significado de nada.
\end{decision}

\subsection{Floats (\emph{carry-over} de la Etapa 1)}

\begin{producciones}
\regla{TipoPrimitivo}{\tl{boolean} \alt{} \tl{char} \alt{} \tl{int} \alt{} \tl{float}}
\regla{Primitivo}{\tl{true} \alt{} \tl{false} \alt{} \tl{intLiteral} \alt{} \tl{charLiteral} \alt{} \tl{floatLiteral} \alt{} \tl{null}}
\end{producciones}
Dos alternativas nuevas, cero conflictos. Sin ellas, el logro de floats de la Etapa 1 queda
inutilizable en la Etapa 2: el léxico emite \tl{float} y \tl{floatLiteral} y el analizador
sintáctico los rechazaría.

\subsection{Impacto sobre el analizador léxico}

Cuatro de los logros exigen elementos léxicos que \textbf{no figuran en \emph{Aspectos Léxicos de
\textsc{MiniJava}}}: la lista de palabras clave de \S2.1 llega hasta \tl{false} y la de operadores
de \S2.4.2 hasta \tl{++} y \tl{-{}-}.

\begin{center}
\begin{tabular}{@{}lll@{}}
\toprule
\textbf{Elemento} & \textbf{Logro} & \textbf{Estado en la Etapa 1} \\
\midrule
\tl{for}     & Fors! E2                 & se reconoce como \tl{idMetVar} \\
\tl{private} & Visibilidad Mejorada E2  & se reconoce como \tl{idMetVar} \\
\tl{?}       & Operador Ternario E2     & es un error léxico \\
\tl{->}      & MiniLambdas E2           & se reconoce como \tl{-} seguido de \tl{>} \\
\bottomrule
\end{tabular}
\end{center}

Se incorporan al analizador léxico \textbf{condicionados al interruptor de su logro}: con el logro
apagado, \tl{?} vuelve a ser un error léxico y \tl{for} y \tl{private} vuelven a ser
identificadores, exactamente como especifica la cátedra. De ese modo el comportamiento léxico base
de la Etapa 1 se conserva íntegro y verificable, que es lo que exige la condición de atomicidad de
la sección~\ref{sec:atomicidad}. Los tokens \tl{:}, \tl{++} y \tl{-{}-}, en cambio, ya existían.
